\section{Discussion}
The experiments reveal a distinction that is hidden by link-level reporting. Increasing the attacker team from one to two does not make self-play unstable; it changes how much of the jammer's marginal disruption the learned radar team can remove. Under the tested task workload, the same radar competence floor coexists with a much larger adversarial loss. The relevant design variable is therefore not only the quality of a single radar policy, but also the number and spatial arrangement of simultaneous interference sources.

The scripted baselines add a second distinction: an equilibrium policy is not a mission-optimal policy. The learned jammer teams cannot punish a greedy mission-stare radar (drop $0.0889$, invariant across all three trained teams), while they punish sweeps ($0.5294$) and the learned radar ($0.3366$). Both sides of the self-play game are specialized to the training opponent distribution: the radar gives up mission throughput to hedge against its co-trained jammer, and the jammer's disruption is effective precisely against attention-spreading radars. Head-to-head drop and containment $\eta$ must therefore be read as properties of the co-adapted pair, not as absolute radar or jammer ratings. This is the same lesson as the singleton-vulnerability result, seen from the other side of the game: static evaluation against out-of-distribution opponents---even non-learning ones---is the cheapest robustness probe available, and it should accompany any self-play claim.

The mechanism control refines this statement. Two jammers at one bearing already reduce containment from $62.8\%$ to $26.6\%\pm2.9\%$ across three control seeds. Cross-fire reduces it further to $24.2\%$. The benchmark therefore assigns primary responsibility to attacker count and a secondary penalty to geometry. This decomposition is more useful than attributing the entire effect to a visually appealing cross-fire explanation: a defense that survives one jammer should be stress-tested against the energy and timing consequences of two sources even when their directions are similar.

R5-lite suggests a practical training response. Mixing opponent classes reduces singleton relative leverage, and the gain saturates near equal exposure to pair and singleton opponents. The result motivates opponent-distribution design as a complement to longer self-play. It does not imply that a full AlphaStar-style league is required for the S7 equilibrium itself \cite{vinyals2019alphastar}. In this benchmark, single-policy self-play converges after approximately 1700 iterations. A population or league becomes relevant when the design objective includes robustness to threat compositions that were not present in the training game---and the counter-adaptation control makes the point concretely: jammers co-trained against the greedy stare recover part of the penalty that self-play never discovered (Section~V-C), so training-time opponent coverage, not longer optimization, is what closes exploit gaps.

The pre-training contestability gate provides a cheap way to decide whether a proposed physical regime can support such a study. It tests whether the best radar response spans a useful probability interval before expensive optimization begins. In S7, the gate predicted that cross-fire would reduce selected sectors while keeping the game nontrivial; the trained behavior and co-located control followed that direction. This gate is a design diagnostic, not a substitute for training or a proof of equilibrium.

The physics contract is deliberately minimal, and the direction of each simplification can be stated. Broadband noise jamming with rectangular spectra is the canonical denial waveform; smarter coherently-informed jamming would only strengthen the attacker, so the attacker-side conclusion is conservative, as is the incoherent cross-jammer power summation argued in Section~III. Idealized phase-aligned cell aggregation within one jammer slightly favors the jammer over a calibrated real array. The exclusions cut in the opposite direction: real radars retain Doppler/MTI filtering and waveform agility that our radar agents do not exercise, so the defense as modeled is weaker than a fielded system, and the h2h values measure scheduling policy under the stated budget rather than absolute electronic-protection capability. The coarse $5\times5$ beam grid quantizes pointing symmetrically for both sides. Within this contract the conclusions are exact. Regime dependence was tested both ways: re-evaluating the converged teams at shifted SNR$_0$ (policies trained at 12~dB) gives $\eta=7.6\%$/$12.0\%$ at 9~dB and $28.0\%$/$39.5\%$ at 15~dB (cross-fire/co-located), and retraining from scratch at the shifted regimes reproduces the same picture---$\eta=2.1\%$ at 9~dB and $23.0\%$ at 15~dB versus $20.2\%$ at 12~dB. Containment is monotone in SNR$_0$ under both protocols, the cross-fire-below-co-located ordering is preserved at every regime, and the 9-dB collapse persists under retraining, identifying the baseline detection margin---not policy adaptation---as the binding constraint there. Geometry and waveform-family breadth remain open.

\section{Limitations}
Each limitation below is paired with the extension that would test it; none invalidates the protocol within its stated scope. First, the physical model uses a single five-sector azimuth grid, two services, fixed site bearings, and a simplified broadband-jammer link budget. The count--geometry decomposition therefore applies to the tested geometry family; a grid over site placements and SNR regimes, with the same gate applied to each configuration, is the corresponding extension, and the retrained 9/15-dB points already cover the regime axis. The greedy-stare result is a property of this link budget: a stronger jammer or a smaller pointing margin could close the stare gap (Section~V-C gives the analytic condition), and the counter-adaptation control confirms the margin is learnable to attack (drop $0.213$ at 2000 iterations, still trending upward). Second, the study is simulation-only. Mission drop is a model-level operational endpoint and does not establish fielded detection, tracking, or electronic-support performance; hardware-in-the-loop replication of the four-view protocol is the natural next tier. Third, the main S7 equilibrium, the S6 baseline, and the co-located mechanism control each have three training seeds; R5-lite uses one seed per mixture condition. Fourth, R5-lite skips the jammer update on singleton iterations, so its absolute drop values confound opponent exposure with jammer gradient budget. The ratio-level result is the intended readout, while gradient-aligned multi-seed mixing remains necessary for a stronger causal intervention claim.

\section{Conclusion}
We presented a task-level benchmark for adversarial co-adaptation between multifunction radars and cooperative jammers. The benchmark links array-factor physics and finite energy to bearing-specific missions with deadlines, then evaluates both sides with controls that separate natural scheduling loss from adversarial loss.

Across three converged seeds, adding a second jammer reduced learned-radar containment from $62.8\%\pm1.6\%$ in the valid three-seed S6 baseline to $23.0\%\pm1.1\%$ at the same S7 team budget, while idle-jammer radar loss remained low. A three-seed co-located control showed that two jammers under the separated S7 radar geometry already produce most of the reduction, with cross-fire adding a smaller penalty. Longer continuation training confirmed a stationary high-loss equilibrium rather than a non-transitive training cycle. Finally, a lightweight opponent-class mixture reduced singleton relative leverage with diminishing returns near 50\% exposure.

Three engineering insights transfer beyond this specific benchmark. First, attacker count is the primary stress variable, which converts directly into a screening rule: a learned scheduler that passes single-jammer validation must be re-tested against two sources---including two co-located sources, which reproduce most of the containment loss---before advancing to the next development stage. Second, an equilibrium policy is not a mission-optimal policy: learned schedulers hedge against their training opponent at the cost of task throughput, so any self-play claim should be accompanied by cheap static evaluations against out-of-distribution opponents---including non-learning heuristics, which exposed the jammer side here as clearly as the singleton exposed the radar side. Third, pointing discipline is the strongest single lever in the modeled budget: mission-stare margins survived every trained jammer team, which suggests that guaranteeing a per-mission pointing margin deserves at least as much design attention as learning the scheduler built on top of it. The next rigorous extensions are a gradient-aligned, multi-seed mixture study, a broader geometry and SNR family, and counter-adaptation training against the stare heuristic, followed by hardware-in-the-loop validation.

\section*{Reproducibility}
The complete evaluation pipeline is deterministic and self-contained: fixed train and validation scenario manifests, fixed action seeds, atomic checkpoints, 18 automated environment and physics gates, and a results-integrity script that re-derives every number quoted in this paper from the raw evaluation JSON files and fails on any mismatch. Upon acceptance, the authors will deposit as supplementary material: (i) the environment, trainer, and evaluation source code; (ii) the training and validation manifests; (iii) all final-evaluation and baseline JSON results for every reported seed; (iv) the figure-generation scripts and the single-source results table; and (v) the three converged checkpoints, so that every table and figure in the paper can be regenerated with one command per figure. A public repository with a citable DOI (IEEE DataPort or equivalent) will accompany the final version; the frozen internal revision identifier is retained in the project log until then.
